\chapter{Conclusion}
\label{ch:conclusion}

This thesis investigated the integration of Retrieval-Augmented Generation (RAG) systems into on-premise Product Lifecycle Management environments through the design, implementation, and evaluation of an AI-powered documentation assistant for Sovelia Core PLM platform. The research addressed the knowledge accessibility challenges characteristic of enterprise PLM deployments, where fragmented documentation spanning vendor help files, release notes, and customer-specific standard operating procedures impedes user productivity and increases support burden, as detailed in Chapter~\ref{ch:introduction}. Operating under strict data governance constraints inherent to manufacturing organizations, this study employed Design Science Research methodology combined with a single case study approach (Chapter~\ref{ch:research-methodology}) to develop an artifact that balances technical sophistication with practical deployability in enterprise contexts.

\section{Research contributions}

This research provides both theoretical and practical contributions to the understanding and implementation of RAG-based documentation assistance in on-premise enterprise environments.

\subsection{Theoretical contributions}

From a theoretical perspective, this thesis extends knowledge about RAG system design in enterprise PLM contexts through several key findings. The research demonstrates that documentation quality, rather than retrieval algorithm sophistication, constitutes the primary determinant of response quality in domain-specific RAG deployments, as evidenced by the evaluation results presented in Chapter~\ref{ch:evaluation-and-results}. This finding challenges the assumption that increasingly complex retrieval strategies necessarily improve system effectiveness, revealing instead that simple vector similarity search with BM25 reranking performs adequately when documentation coverage is comprehensive and well-structured.

The study provides empirical evidence that hybrid cloud-local architectures offer viable deployment paths for organizations balancing data governance constraints with performance requirements. By demonstrating that embedding generation can occur on-premise while LLM inference leverages cloud APIs (Chapter~\ref{ch:system-design-and-implementation}), the research identifies a practical middle ground between complete on-premise deployment (with significant infrastructure requirements) and full cloud migration (often prohibited by data sovereignty constraints). This architectural pattern addresses a gap in existing literature (Chapter~\ref{ch:literature-review}), which predominantly assumes either fully cloud-based or fully on-premise deployments without exploring hybrid approaches.

The evaluation methodology contributes a factored assessment framework specifically adapted for RAG systems in technical documentation contexts. By decomposing response quality into four independent dimensions (factual accuracy, relevance, source attribution, completeness) assessed through expert evaluation, the study provides a replicable evaluation approach that enables precise identification of system strengths and weaknesses. This multi-dimensional framework extends standard information retrieval metrics to address the unique challenges of RAG-based question answering where both retrieval quality and generation accuracy must be evaluated.

The research also establishes empirical foundations for understanding performance trade-offs between different LLM deployment configurations. The performance comparison between cloud-based OpenAI GPT-5 Nano (25,9 seconds) and on-premise Ollama GPT-OSS-20B on dedicated GPU infrastructure (91,4 seconds with 4$\times$ NVIDIA L40S, 192 GB VRAM total) reveals that inference framework optimization and architectural efficiency can substantially impact response times beyond raw computational capacity. This specific case demonstrates that Ollama's single-GPU architectural constraints and optimization level create performance gaps that more sophisticated frameworks (such as vLLM with multi-GPU parallelization) could substantially reduce. These findings inform deployment decisions for organizations evaluating infrastructure investments versus operational costs, while highlighting that framework and optimization choices significantly influence on-premise deployment viability.

\subsection{Practical contributions}

From a practical perspective, this thesis delivers a production-ready RAG system integrated with Sovelia Core that demonstrates operational viability through expert evaluation and performance testing (Chapter~\ref{ch:evaluation-and-results}). The implementation achieves mean expert evaluation scores between 4,31 and 4,49 on a 5-point scale across factual accuracy, relevance, source attribution, and completeness dimensions, corresponding to quality ratings between ``Good'' and ``Excellent.'' This performance level indicates that the system provides sufficient value for deployment in operational PLM environments, meaningfully reducing documentation search time and improving information accessibility for users.

The system architecture provides a reusable reference implementation addressing common challenges in on-premise RAG deployment (Chapter~\ref{ch:system-design-and-implementation}). Key architectural elements include database-level access control integration maintaining PLM permission boundaries, provider abstraction layers enabling flexible LLM selection without vendor lock-in, agentic tool-calling patterns allowing multi-step retrieval when necessary, and documentation pipeline automation supporting periodic updates from multiple source types. These components collectively demonstrate how RAG capabilities can be integrated into existing enterprise systems without requiring extensive infrastructure changes or fundamental architectural reorganization.

The documentation pipeline implementation establishes operational patterns for managing RAG system lifecycle in production environments. The pipeline processes documentation from three distinct sources (internal PLM documents, vendor help files via Document360 API \parencite{Document360API}, and release notes from web content), performing chunking, embedding generation, and version-aware updates to maintain current knowledge bases. This multi-source integration pattern addresses a practical challenge frequently encountered but rarely documented in RAG literature: how to systematically maintain documentation currency as sources evolve independently.

The cost analysis provides decision-making frameworks for organizations evaluating deployment options. With cloud-based LLM inference costs of approximately \$0,001 per query (\$30 monthly for 1 000 daily queries), the research demonstrates that API expenses remain economically justifiable for interactive assistance scenarios where sub-30-second response times enhance user experience. This quantification enables direct comparison against infrastructure investments required for equivalent on-premise performance, informing build-versus-buy decisions for enterprise AI capabilities.

\section{Addressing the research questions}

The research questions posed in Chapter~\ref{ch:introduction} guided the investigation throughout this study. The following sections summarize how each question was addressed.

\subsection{RQ1: RAG system architecture for on-premise PLM environments}

The first research question asked: \textit{What RAG system architecture best fits Sovelia Core's on-premise, PostgreSQL-based environment and data governance constraints?}

The implemented architecture (detailed in Chapter~\ref{ch:system-design-and-implementation}) demonstrates that effective RAG systems can be built using existing PostgreSQL infrastructure augmented with the pgvector extension, avoiding the need for specialized vector databases. This approach leverages familiar database technologies, reduces operational complexity, and maintains unified data management for both structured PLM data and vector embeddings. The five-layer architecture (data, embedding, business logic, API, presentation) provides clear separation of concerns while maintaining integration with existing PLM access control mechanisms.

The hybrid deployment model, where embedding generation occurs on-premise while LLM inference optionally utilizes cloud APIs, addresses data governance constraints by ensuring that documentation content remains within controlled infrastructure during retrieval. This architecture enables organizations to balance performance requirements (leveraging optimized cloud inference) with data sovereignty concerns (maintaining control over sensitive information). For organizations requiring complete on-premise operation, the provider abstraction layer supports local LLM deployment via Ollama, though with documented performance trade-offs.

Database-level access control integration ensures that the RAG system respects existing permission boundaries, preventing information leakage through AI responses. This design treats the AI assistant as an extension of the PLM interface rather than a separate system with independent permissions, maintaining security models that organizations have already validated and audited.

\subsection{RQ2: Documentation pipeline engineering}

The second research question asked: \textit{How should customer documentation pipelines (parse $ \rightarrow $ chunk $ \rightarrow $ embed $ \rightarrow $ retire) be engineered for reliability, maintainability, and cost efficiency?}

The implemented documentation pipeline demonstrates reliable multi-source processing through standardized extraction interfaces for internal PLM documents, external API-accessed vendor documentation, and web-scraped content. Version-aware processing prevents redundant re-embedding of unchanged content while ensuring that updates propagate correctly through the system. Cascade deletion behavior automatically cleans up associated chunks when source documents are removed, preventing orphaned data accumulation.

The chunking strategy (1 000-character chunks with 200-character overlap) balances context completeness with retrieval precision, ensuring that individual chunks contain sufficient information to be meaningful while maintaining manageable vector search result sets. This parameter selection aligns with patterns observed in comparable technical documentation RAG implementations reviewed in Chapter~\ref{ch:literature-review} and proved effective across diverse query types during evaluation (Chapter~\ref{ch:evaluation-and-results}).

Cost efficiency is achieved through local embedding generation using the open-source Nomic Embed Text model, eliminating per-embedding API costs while maintaining adequate semantic representation quality. The pipeline architecture supports scheduled batch updates, allowing documentation synchronization during off-peak hours without impacting interactive query performance. Automated scheduling through stored procedures enables hands-off maintenance once initial configuration is established.

Maintainability is addressed through modular pipeline components with clear interfaces, comprehensive error logging to facilitate troubleshooting, and configuration-driven behavior enabling parameter adjustments without code modifications. The provider abstraction pattern ensures that embedding model selection remains flexible, allowing organizations to adopt newer models as they become available without rewriting pipeline logic.

\subsection{RQ3: Risk mitigation in practice}

The third research question asked: \textit{Which risks, particularly those related to access control and stale context, are most critical, and what mitigations work in practice?}

The evaluation (Chapter~\ref{ch:evaluation-and-results}) revealed that documentation coverage gaps represent the most critical risk to system utility, with evaluator feedback consistently attributing response weaknesses to missing or incomplete source material rather than technical failures in retrieval or generation. This finding emphasizes that RAG deployment requires concurrent investment in documentation quality improvement, as sophisticated algorithms cannot compensate for absent information. Practical mitigation involves systematic documentation audits identifying coverage gaps, particularly in integration scenarios, troubleshooting procedures, and configuration options that users frequently query.

Access control risks are mitigated through database-level permission filtering during retrieval, ensuring that unauthorized content never reaches the language model and cannot influence generated responses. This approach implements the principle that access control enforcement should occur during retrieval rather than after generation, preventing information leakage while preserving answer quality. The integration with existing PLM permission mechanisms ensures consistent security models across traditional and AI-powered interfaces.

Stale context risks, where outdated documentation produces incorrect guidance, are addressed through version-aware pipeline processing that tracks document currency and supports periodic updates from authoritative sources. The multi-source architecture maintains distinct update schedules for internal PLM documents (modified through system workflows), vendor help files (synchronized via Document360 API), and release notes (extracted from web content), ensuring that each source remains current according to its natural update cycle.

Hallucination risks, where the language model generates plausible but incorrect information, are mitigated through source attribution requirements and retrieval-focused prompt engineering. Every generated response includes citations to source documentation, enabling users to verify claims and identify potential fabrications. The agentic architecture explicitly instructs the model to base responses on retrieved context rather than parametric knowledge, reducing reliance on potentially outdated training data.

The evaluation revealed category-specific performance patterns (discussed in detail in Chapter~\ref{ch:evaluation-and-results}) where integration questions (spanning CAD, PLM, ERP system boundaries) exhibited weaker performance (mean scores 3,89--4,22) compared to version-specific questions (4,56--4,78). This pattern indicates that cross-system workflow documentation requires specialized attention, potentially through enhanced retrieval strategies or structured metadata linking related concepts across system boundaries. Recognizing these category-specific risks enables targeted mitigation efforts where generic approaches prove insufficient.

\section{Limitations and future research directions}

Several limitations constrain the generalizability of these findings and suggest promising directions for future research.

\subsection{Study limitations}

The evaluation (Chapter~\ref{ch:evaluation-and-results}) was conducted on a relatively small documentation corpus (113 documents, 217 chunks) representing Sovelia Core's current external documentation. Production PLM deployments often encompass substantially larger knowledge bases including customer-specific customization documentation, legacy system manuals, and accumulated tribal knowledge. Retrieval performance and answer quality may degrade at larger scales, though the sub-100-millisecond retrieval times observed suggest headroom for growth before optimization becomes necessary.

Performance testing occurred on a specific development workstation and GPU infrastructure configuration, with results reflecting this particular hardware setup. Production environments with different specifications, concurrent user loads, or network characteristics may exhibit different performance profiles. The Ollama performance results specifically reflect single-GPU architectural limitations and might improve substantially with alternative inference frameworks optimized for multi-GPU deployments.

Expert evaluation involved three domain experts, providing valuable qualitative insights but limited statistical power for detecting subtle quality differences. The systematic variation in absolute rating levels between evaluators (product manager consistently 0,6--0,8 points lower than product expert) highlights calibration differences that larger evaluator panels might mitigate through consensus processes.

The 30-question evaluation set, while carefully constructed to represent diverse query types, may not comprehensively capture the full distribution of information needs in production usage. Certain query patterns (multi-part questions, highly specific customization scenarios, queries requiring temporal reasoning across multiple releases) may be underrepresented. The evaluation also focused on text-based queries, not addressing scenarios involving CAD model references, numerical data, or graphical information.

The study employed expert evaluation rather than end-user testing with actual Sovelia Core users. While domain experts can assess technical accuracy and documentation correctness, they may not fully represent novice user perspectives, information seeking patterns, or tolerance for different error types. User satisfaction, task completion rates, and time-to-resolution metrics would provide complementary insights into practical utility.

\subsection{Future research directions}

Several promising research directions emerge from this study's findings and limitations.

\textbf{Documentation coverage optimization.} Systematic investigation of documentation characteristics that predict RAG system performance could inform targeted improvement efforts. Research questions include: What documentation structures best support retrieval effectiveness? How should procedural information be organized to maximize answer completeness? What metadata schemas enhance cross-system integration query handling? Developing documentation quality metrics correlated with RAG response quality would enable proactive gap identification and prioritization of documentation improvement efforts.

\textbf{Advanced retrieval strategies for integration queries.} The notably weaker performance on integration questions (3,89--4,22) identified in Chapter~\ref{ch:evaluation-and-results} suggests value in specialized retrieval approaches for multi-system workflow scenarios. Potential directions include knowledge graph construction linking CAD, PLM, and ERP entities (similar to approaches reviewed in Chapter~\ref{ch:literature-review}); multi-stage retrieval explicitly targeting system boundary documentation; and hybrid approaches combining structured metadata about interfaces with unstructured documentation retrieval. Comparative evaluation of these strategies on integration-focused query sets would quantify potential improvements and inform deployment decisions for organizations prioritizing cross-system workflow support.

\textbf{On-premise performance optimization.} Investigation of inference frameworks optimized for multi-GPU deployments (e.g., vLLM, TensorRT-LLM) could substantially reduce the performance gap between cloud-based and on-premise LLM deployments. Research comparing inference latency, throughput, and resource utilization across different frameworks and model architectures would inform infrastructure decisions for organizations prioritizing data residency. Quantifying the point at which query volume justifies GPU infrastructure investments versus cloud API costs would provide deployment decision frameworks enabling evidence-based build-versus-buy analyses.

\textbf{Temporal reasoning and version awareness.} The mixed performance on version-specific questions suggests opportunities for enhanced temporal reasoning capabilities. Approaches might include explicit version metadata in chunk embeddings, temporal filtering during retrieval, specialized prompts emphasizing chronological distinctions, or hybrid systems combining structured release databases with unstructured documentation retrieval. Evaluation on version-comparison and feature-evolution queries would assess these enhancements and determine whether temporal information integration justifies implementation complexity.

\textbf{Long-term deployment studies.} Longitudinal evaluation tracking system performance as documentation evolves would address temporal limitations in this study. Research questions include: How does retrieval effectiveness degrade as documentation volume grows? What maintenance processes preserve answer quality as content updates? How do users' information seeking patterns evolve with system familiarity? Production deployment with usage logging and periodic quality assessment would yield insights into operational sustainability and identify early warning indicators of performance degradation requiring intervention.

\textbf{Multi-modal documentation support.} PLM documentation often includes CAD model visualizations, technical diagrams, tables of specifications, and graphical user interface screenshots. Extending the RAG system to handle visual information through multi-modal embeddings, image-text retrieval, and vision-language models would address a significant gap in current capabilities. Evaluation on queries requiring visual information (``show me the location of the BOM tab,'' ``what does the lifecycle state diagram look like'') would quantify multi-modal extensions' utility and determine whether visual information integration provides sufficient value to justify implementation costs.

\textbf{User-facing evaluation and deployment.} Controlled user studies with actual Sovelia Core users would complement expert evaluation with real-world usage insights. Metrics including task completion rates, time-to-resolution, user satisfaction scores, and preference comparisons against traditional documentation search would quantify practical impact on daily workflows. Production deployment with opt-in user testing would enable large-scale behavioral analysis revealing actual information seeking patterns and system utility in authentic work contexts, providing insights that laboratory evaluations cannot capture.

\section{Closing remarks}

The integration of RAG-based AI assistants into on-premise PLM environments represents both a technical achievement and a pragmatic response to persistent knowledge management challenges in enterprise software systems. This thesis demonstrates that effective documentation assistance can be achieved within the constraints characteristic of manufacturing organizations: strict data governance requirements, existing infrastructure investments, and operational complexity concerns.

The evaluation results (Chapter~\ref{ch:evaluation-and-results}) validate the central premise that RAG systems provide operationally viable quality for PLM documentation support, achieving expert evaluation scores between ``Good'' and ``Excellent'' across multiple quality dimensions. These findings suggest that organizations can meaningfully improve user productivity and reduce support burden through AI-powered documentation assistance without requiring fundamental architectural changes or prohibitive infrastructure investments.

Critically, the research reveals that documentation quality, rather than algorithmic sophistication, determines system effectiveness. This finding has important implications for deployment strategy: successful RAG implementation requires concurrent investment in source material improvement alongside technical infrastructure. Organizations pursuing AI-enhanced documentation assistance must recognize that retrieval and generation technologies amplify existing documentation strengths while exposing weaknesses with equal fidelity.

The hybrid deployment architecture demonstrated in this study (Chapter~\ref{ch:system-design-and-implementation}) offers a practical path forward for organizations navigating the tension between performance requirements and data governance constraints. By separating on-premise retrieval from optionally cloud-based generation, the approach enables organizations to leverage optimized inference infrastructure while maintaining control over sensitive information. This architectural pattern merits broader adoption in enterprise AI contexts where similar constraints and objectives prevail.

As RAG technologies continue to evolve and on-premise LLM inference frameworks improve, the performance gap between cloud-based and local deployments will likely narrow. Organizations implementing RAG systems today should prioritize architectural flexibility through provider abstraction patterns, ensuring that infrastructure choices remain reversible as technological and organizational contexts change. The rapid evolution of language models, measured in months rather than years, contrasts sharply with PLM system lifecycles spanning 7-15 years, emphasizing the importance of designs that accommodate continuous adaptation.

This research contributes to the growing body of knowledge about practical AI integration in enterprise environments, demonstrating that sophisticated capabilities can be deployed within realistic organizational constraints when architectural decisions reflect actual deployment contexts. The findings provide both theoretical insights about RAG system design principles and practical guidance for organizations implementing similar systems in comparable environments. As manufacturing organizations increasingly explore AI-powered enhancements to existing systems, evidence-based implementation patterns derived from real-world deployments prove essential for distinguishing viable approaches from speculative proposals.

The continued development and refinement of RAG-based documentation assistance in PLM environments promises substantial benefits for organizations willing to invest in both technical infrastructure and source material quality. This thesis establishes that the foundational capabilities work effectively, providing a platform upon which future enhancements can build toward increasingly sophisticated and valuable AI-powered knowledge management in enterprise contexts.